\documentclass[11pt]{article}
\usepackage{aiguide}
\guidekicker{Guide 1 of 3}
\begin{document}
\makeguidetitle{Studying Mathematics with AI}{A method for using chat assistants without letting them do the learning for you.}

An AI assistant can explain an idea in five different ways, at midnight, without
losing patience. It can also hand you a finished solution that you copy, nod
at, and forget by morning. Used well, it is a tutor. Used badly, it is a way to
skip the part of mathematics that makes you better at it: getting stuck and
finding the way out yourself.

\section{Five rules}
\begin{enumerate}[label=\textbf{\arabic*.}]
  \item \textbf{Ask for hints, not answers.} A finished solution feels like
    progress but teaches little. Ask for the next hint, then do the step
    yourself.
  \item \textbf{Redo every step by hand.} Close the chat and solve the problem
    again on paper. If you cannot, you have not learned it yet.
  \item \textbf{Check every number and every claim.} Substitute your answer
    back into the equation, test a value, or run it through a computer algebra
    system such as Wolfram|Alpha or GeoGebra. Look up any theorem or source a
    model names; they are sometimes invented.
  \item \textbf{Say when you used AI.} Follow your teacher's or your course's
    rules. When in doubt, say what you used it for.
  \item \textbf{Keep personal data out.} Never paste names, grades, photos,
    passwords, or anything about classmates into a chat.
\end{enumerate}

\section{A study session in three steps}
\begin{description}[leftmargin=0pt,labelsep=0.4em]
  \item[Before.] Try the problem alone for at least ten minutes. Write down
    exactly where you got stuck: that sentence is your question.
  \item[During.] Use one of the prompts below. Answer the tutor's questions
    in your own words. If it gives away the solution, say so and ask it to go
    back to hints.
  \item[After.] Solve a similar problem with the chat closed. Check the result
    yourself (Guide~3 lists what to look for). Note the idea you were missing.
\end{description}

\begin{notebox}
\textbf{Not during tests or graded work,} unless your teacher or course says
you may. Practice is where AI helps; an assessment checks what \emph{you} can do.
\end{notebox}

\section{Prompts to copy}
Replace the parts in brackets before you send them. They work in ChatGPT,
Gemini, Claude, and similar assistants. ChatGPT's study mode and Gemini's
Guided Learning are built-in versions of the first prompt.

\begin{promptbox}{Socratic tutor}
Act as a Socratic tutor. I am working on this problem: [paste the problem]. Do
not give me the solution. Ask me one question at a time to help me find the
next step, and wait for my answer before you continue.
\end{promptbox}

\begin{promptbox}{Check my solution without solving it}
Here is a problem and my solution: [paste both]. Do not write a full solution.
Tell me whether each step is correct. For the first wrong step, say which line
it is and give me a hint, not the fix.
\end{promptbox}

\begin{promptbox}{Quiz me}
Quiz me on [topic] at the level of [class or course]. Ask five questions, one at
a time, from easy to hard. After each answer, tell me whether I am right and
why, then ask the next one.
\end{promptbox}

\begin{promptbox}{Explain it at my level}
Explain [concept] to a student in [class or course]. Use one short worked
example with small numbers, then give me a similar exercise to try on my own. Do
not solve the exercise.
\end{promptbox}

\begin{promptbox}{Find the error in this step}
This step is wrong, but I cannot see why: [paste the line before and the line
after]. Point to the mistake and name the rule I broke. Do not continue the
solution.
\end{promptbox}

\begin{promptbox}{Make a study plan}
I have a test on [topics] in [number] days and can study [time] per day. Make a
day-by-day plan that mixes review, practice problems, and one practice test.
Keep the last day for revision only.
\end{promptbox}

\section{Why this works}
You remember what you produce, not what you read. A hint leaves the hard step
to you; a full solution takes it away. Checking the answer yourself is also
practice: it is the same skill you need when no one, human or machine, is there
to check it for you.

\end{document}
